\begin{introduction}
\label{chapitre:introduction}

Les technologies de séquençage d'ADN permettent aujourd'hui d'analyser de grandes quantités de
matériel génétique. Alors que la première génération de technique de séquençage (à partir des années
1970) était limitée à un faible débit, les technologies de deuxième génération (aussi appelées
<<nouvelle génération>> ou \textit{NGS}, à partir des années 2000) ont grandement augmenté la
vitesse de séquençage mais demeuraient limitées à de courtes séquences d'au plus quelques centaines
de nucléotides. Les technologies de troisième génération (à partir des années 2010) constituent une
autre amélioration qui permet d'analyser des séquences encore plus longues, jusqu'à plusieurs
milliers de nucléotides, encore plus rapidement \cite{satam2023}. Ces technologies ont contribué à
améliorer nos connaissances en génomique et mené à des applications en biologie et médecine, par
exemple, pour la détection de pathogènes, les diagnostics cliniques, la protection de
l'environnement et l'agriculture \cite{liu2025}. De nombreux outils ont aussi été développés pour
analyser les données de séquençage et faciliter le travail des chercheurs \cite{roy2024}. Ce mémoire
s'intéresse en particulier à une application de l'analyse de séquences d'ADN, la classification
taxonomique de métagénomes, qui vise à identifier les espèces présentes dans des microbiomes.

La métagénomique se heurte à des problèmes majeurs. Le \textbf{volume de données} complique les
analyses. Un séquençage conventionnel de métagénome avec des technologies de troisième génération
peut générer $10\;Go$ de données par expérience \cite{PerezCobas2020}. Les bases de données
contiennent aussi beaucoup d'information : RefSeq \cite{refseq} contient des génomes de référence de
haute qualité pour plus de $176\;395$ organismes en avril 2026 et GTDB \cite{gtdb} contient un
taxonomie exhaustive pour plus de $153\;000$ espèces de procaryotes. Identifier une espèce en
comparant une séquence d'ADN avec une base de données de référence entraînerait un temps de calcul
inacceptable pour des applications réelles, c'est pourquoi de nombreuses méthodes d'analyse ont été
développées pour obtenir des résultats plus rapidement \cite{roy2024}. La \textbf{fiabilité des
résultats} fait référence à la confiance qu'on peut accorder aux données produites par une analyse
métagénomique. Gihawi et al expliquent qu'une étude métagénomique influente était basée sur une
analyse erronée : en normalisant mal les données et en utilisant des séquences contaminées par l'ADN
humain, les chercheurs de cette étude en étaient venus à la conclusion qu'il est possible de prédire
certains types de cancers à partir de données de microbiotes humains, ce qui a induit d'autres
chercheurs en erreur \cite{gihawi2023}. Il est donc impératif de développer des techniques d'analyse
qui soient non seulement efficaces, mais aussi rigoureuses et facile à interpréter pour éviter
d'autres scénarios du même genre.

Ce mémoire vise à améliorer des techniques de classification taxonomique de données de séquençage
métagénomique, ce qui implique les étapes principales suivantes :

La \textbf{génération de données synthétiques} permet de créer des ensembles de données pour
entraîner et valider les outils de classification. Les données doivent être générées à partir de
séquences de haute qualité et représenter des échantillons métagénomiques variés pour éviter de
surajuster les modèles et d'obtenir des performances non représentatives d'un cas d'utilisation
réel.

L'élaboration de \textbf{modèles de classification taxonomique} consiste à concevoir des réseaux
neuronaux inédits pour améliorer les techniques existantes. Ces modèles doivent offrir des
performances supérieures en terme d'utilisation de ressources informatiques (temps de calcul et
utilisation de la mémoire) et de classification (précision plus élevée).

L'\textbf{entraînement des modèles} consiste à modifier automatiquement les paramètres des modèles
pour leur permettre d'effectuer les classifications taxonomiques. Cette étape requiert des
ressources intensive et doit s'appuyer sur de bons hyperparamètres pour converger vers une solution
efficace. Un entraînement mal configuré peut faire échouer le modèle à apprendre comment résoudre
une tâche et ainsi causer un gaspillage de ressources \cite{liu-etal-2020-understanding}.

Ce projet consiste à répertorier des architectures de réseaux neuronaux utilisés en métagénomique et
encore non appliqués à la métagénomique qui ont le potentiel d'offrir de meilleures performances.
Après la sélection des modèles, le processus d'entraînement est ajusté pour chaque modèle afin de
maximiser leurs performances. Les hyperparamètres sont modifiés et différentes techniques
d'entraînement sont testées pour déterminer les meilleures manières d'obtenir une classification de
haute qualité. Tous les modèles sont ensuite testés avec les mêmes données de test et d'entraînement
pour identifier les modèles les plus efficaces. On compare aussi les nouveaux modèles développés
dans le cadre de ce projet avec des techniques de classification développées par d'autres équipes de
chercheurs pour confirmer que les nouveaux modèles constituent bel et bien des amélioration par
rapport aux outils déjà existants.

Ce mémoire présente les contributions principales suivantes :

\begin{list}{$\bullet$}{}
    \item Comparaison systématique de différentes architectures de réseaux neuronaux pour la
        classification taxonomique de données métagénomique, en particulier une comparaison
        approfondie entre les transformateurs et les modèles à espace d'état (SSM).
    \item Étude des effets du pré-entraînement et de l'ajustement sur les performances des modèles
        de classification.
    \item Étude des effets du prétraitement des données de séquençage (jetonisation et traitement
        par CNN) sur les performances des modèles de classification.
    \item Conception de modèles de réseaux neuronaux inédits spécifiquement développés pour la
        classification taxonomique.
    \item Analyse de la consommation de ressources informatique par les modèles d'AA.
    \item Publication du code source ouvert pour répliquer tous les résultats du projet.
\end{list}

Ce mémoire comporte quatre sections principales. Le chapitre 1 présente un état de l'art des outils
d'analyse de données métagénomique ainsi que de l'apprentissage profond, avec un accent particulier
sur les réseaux neuronaux dédiés aux analyses de données génomiques. Le chapitre 2 explique la
méthodologie du projet : collecte et transformation des données, conception des modèles,
entraînement et évaluation. Le chapitre 3 présente les résultats des modèles. Le chapitre 4 discute
des résultats et propose des améliorations possibles aux nouveaux modèles.

\end{introduction}
